\clearpage
\setcounter{page}{1}
\maketitlesupplementary
\setcounter{section}{0}
\setcounter{table}{0}
\setcounter{figure}{0}
\setcounter{equation}{0}
\setcounter{algorithm}{0}
\renewcommand{\thesection}{S\arabic{section}}
\renewcommand{\thetable}{S\arabic{table}}
\renewcommand{\thefigure}{S\arabic{figure}}

This supplement complements the main paper with the details needed to assess and reproduce its results. Sec.~\ref{supp:roles} describes the cohort and data roles and our search for an untouched external cohort. Sec.~\ref{supp:construction} gives the implementation checks of masking, the architectures, training and the controls of the compact head. Sec.~\ref{supp:locked} details the evaluation protocol, the uncertainty estimates and the comparisons on top of the full-field ranker. Secs.~\ref{supp:new}--\ref{supp:repro} cover the follow-up comparisons, the anatomy groups and earlier results on the same benchmark, and Sec.~\ref{supp:sizeaudit} reports an exploratory audit by lesion size and proposal displacement. Sec.~\ref{supp:discussion} extends the discussion of the main paper, and Sec.~\ref{supp:tables} collects all detailed results tables. Seeds~1 and~2 of the main paper correspond to training seeds 3691 and 9143.

\section{Cohort, roles and external eligibility}
\label{supp:roles}
\begin{table}[t]
\centering
\small
\caption{Patient-disjoint data roles. The two excluded scans belong to one repeated patient.}
\label{tab:roles}
\begin{tabular}{@{}lrr@{}}
\toprule
Role & Scans & Proposals \\
\midrule
Weights & 356 & 302,680 \\
Epoch selection & 44 & 35,529 \\
Readout development & 45 & 40,162 \\
Cutoff calibration & 88 & 73,542 \\
Evaluation subsets 6--7 & 177 & 150,894 \\
Evaluation subsets 8--9 & 176 & 150,412 \\
Excluded (repeated patient) & 2 & 1,756 \\
\bottomrule
\end{tabular}
\end{table}


We use the LUNA16 false-positive reduction release, which provides the V2 proposals and the reference files~\cite{setio2017luna16}. Its included nodules measure at least $3\mm$ and were accepted by at least three readers under the benchmark definition~\cite{setio2017luna16,armato2011lidc}. Every released proposal is kept at its physical coordinates. Data roles are assigned per patient using the LIDC patient crosswalk (Tab.~\ref{tab:roles}). The two scans of a repeated patient in subsets 1 and 6 are excluded from training and evaluation, and every remaining scan belongs to a unique patient. The evaluation pool combines subsets 6--7 and 8--9 and retains all lesions, including those that no proposal matches. Earlier crop experiments in related work used all nodule-bearing scans of this pool, and prior work on LUNA16 and LUNA25 informed our hypothesis. The pool is therefore reused; patient-disjoint roles and choices fixed before evaluation protect the comparisons between the models reported here.

\paragraph{Search for an untouched cohort.}
We sought an external cohort with scan-level coverage, usable references of included and excluded findings, and a proposal protocol that could be fixed without consulting outcomes. LNDb has 236 public training scans and 58 reserved test scans~\cite{pedrosa2021lndb}, and its reserved image archives can be read. However, the public coordinates of the reserved scans are intended for segmentation and mix target findings with decoys, while the complete detection labels are withheld. Earlier related work used the 177 public training scans with multiple reads, and the other 59 public scans have a single reader, below the reference criterion used here.

LIDC-annot-NLST501 provides four reads per scan and a corrected series crosswalk~\cite{nlst501reference}. We reconciled all 501 subject and series keys with the available exposure inventories. The public LUNA25 roster overlaps with 112 of these patients, including 55 identical original series. A further 71 patients are excluded because they were used to validate earlier detectors (140 such patients in the full 501-case roster before the LUNA25 exclusion). The remaining 318 cases lack a verified crosswalk to the separate hidden LUNA25 tuning partition~\cite{peeters2025protocol}, so we cannot verify that they are untouched.

The Duke screening dataset uses report-informed annotation, assisted by a detector with manual review and additions~\cite{wang2025duke}. Its file listings required access, and the local inventories could not establish a complete detection reference or an eligible untouched roster. Lung TIME mostly contains metastasis workups of adolescents and nine adult scans acquired during treatment, and we could not verify a usable full reference and archive for this comparison~\cite{dolejsi2009time}. The ANODE09 challenge withholds the reference needed for an analysis over all proposals. We therefore used no external outcomes; instead, the single-stream 3D experiment varies the architecture on the original proposals and patients.

\section{Construction and training details}
\label{supp:construction}

\paragraph{Implementation checks of masking.}
For a trilinear sample point $q$, every native voxel $v$ with positive weight satisfies $\lVert v-q\rVert\leq d$, where $d$ is the diagonal of a native voxel. Because the native direction matrix is orthogonal, this physical bound holds in scanner coordinates. The samples at $r_N$ and $r_F$ therefore read only voxels at distances in $[r_N-d,r_F+d]$ from the proposal, and since $r_N-d\geq10.05\mm$, these sources are disjoint from the native ball. Sample points outside the scan read a fixed value of air. The implementation masks a neighbourhood of the proposal with a margin of native voxels for interpolation, and the model receives clipped and quantized intensities without a separate visibility channel.

We compared this implementation with literal masking of the whole volume in nine rotated, reflected, anisotropic and boundary configurations. Randomizing 485,801 hidden values left all crops, samples and background levels unchanged, and the check covered 15,664 interpolation contributions with positive weight. The implementation of the native-source condition passed nine equivalent checks against literal masking and against randomized hidden sources. For every released proposal, it reproduced the sample values before the crop was quantized and recomputed the median background level from the corresponding subset of samples. A retained source enters the centred crop if, along every axis, some in-bounds grid point lies strictly less than one native voxel step away; this follows from the separability and positivity of the trilinear weights, and we checked the resulting counts against an explicit enumeration of grid neighbours on the test configurations.

\paragraph{Architectures.}
The cached crop spans $37.2\mm$ between the first and the last of 32 sample centres along each axis, and inference uses the centred $28^3$ subgrid. The two-scale model processes a central $16^3$ view and a $14^3$ average-pooled context view. Each stream has a $3^3$ stem with 16 channels followed by three stride-two stages of widths 32, 64 and 96. Mean and max statistics feed a 96-unit projection with layer normalization, SiLU and dropout 0.25, giving an image ranker with 1,283,873 parameters. The three-view 2D model takes six central slices per orientation, encodes them with a shared 2D encoder and concatenates the three mean/max summaries.

The single-stream 3D model rearranges each non-overlapping $2^3$ neighbourhood into eight channels on the $14^3$ grid, without averaging. A 32-channel convolutional stem is followed by residual blocks of widths 32, 64 and 96 with strides 1, 2 and 2; each block uses two $3^3$ convolutions and a projection shortcut when the dimensions change. Pooled mean and max features of the last two stages feed a 96-unit projection and a scalar classifier, for a total of 683,265 parameters. The single-stream and native-source models use seed~1 only.

\paragraph{Training and selection.}
All image rankers are trained for 48 epochs of 160 batches, each with 24 positives and 104 negatives. Labels follow strict reference matches, and ambiguous benchmark positives and excluded findings are not sampled as negatives. Each training crop is shifted by a uniform discrete offset of zero to four grid steps from the cache corner, rotated by in-plane right angles, flipped independently along each axis, and perturbed by intensity gain and offset jitter; the visibility mask is applied before these transformations. We optimize with AdamW at a base learning rate of 0.0018, with cosine decay to 0.12 of this rate, weight decay 0.012 and bfloat16 computation on the GPU. Every fourth epoch, we score all proposals of the checkpoint-selection role and keep the checkpoint with the highest capped CPM, breaking ties by the earliest epoch. At each of these epochs except the last, we refresh a pool of the 16,000 highest-scoring training negatives, from which subsequent batches draw half of their negatives.

The sparse readout heads receive the unaugmented centred logit and the unaugmented sample values. They are trained for 24 epochs with the same batch size and schedule shape, using a hard-negative pool derived once per seed from the frozen ball-only ranker, and each head is selected by CPM on the checkpoint-selection role every four epochs, again breaking ties by the earliest epoch. The direction-blind and direction-aware heads of width 48 have 53,977 parameters each and those of width 96 have 212,785 each; the compact head with 873 parameters is refitted on the same batches. The attention blocks use pre-normalization, four heads and a two-layer feedforward network with expansion factor two, and symmetric mean and max pooling summarize the final normalized tokens. A width-matched projection of the clipped logit $i$ and the exterior summary $e$ enter an MLP on $[i,e,i\odot e]$, which predicts a scalar residual; its zero-initialized last layer makes each head start from the logit of the ball-only ranker. The token permutation moves both axes of the directional bias matrix together with the observations. A test fixture confirms that permuting the input and the bias jointly preserves the spatial encoding, whereas permuting the input alone can change it; the direction-blind encoding is invariant to either permutation, and both heads are invariant to swapping the two half-axes of an axis.

\paragraph{Initial compact heads.}
Let $s$ denote the logit of the image ranker and $z$ its penultimate embedding. The logit pathway uses the baseline $B_L=s$, and the embedding pathway uses a separately trained image-only score $B_E=s+q(z)$; a scalar or endpoint residual is added to the respective baseline. Every variant uses the same 8-wide projections of the interior and the exterior and a 16-wide SiLU MLP on the two projections and their elementwise product, and the scalar and morphology variants represent their single value as a one-row set. The paired, unpaired and shuffled-pair variants use tensors of identical size, and the embedding residual has 1,633 parameters. Gradients stop at the detached image features and the corresponding image-only score. These heads were trained alongside the image ranker, using its checkpoint and its mining of hard negatives, with the samples kept at the original proposal; the refits described above keep both modalities at the same centred, unaugmented pose.

\paragraph{Pairing controls.}
The shuffled-pair control reassigns whole half-axes between opposite directions while keeping both radii within each half-axis. A fixed bank of 512 bijections enforces at least $45^\circ$ of undirected separation between the source and the destination axis, and neither the axis identity nor the chosen permutation is given to the network. The unpaired control pairs each half-axis with itself. In a synthetic test with identical scalar values and half-axis marginals, but distinct directional clusters in fitting and evaluation, the paired head reaches 100\% accuracy on 320 held-out examples, while each of eight shuffled-pair draws and the unpaired head reach 50\%. This test verifies that pairing information is available to the heads; it provides no evidence about patient-level performance.

\section{Evaluation protocol, uncertainty and operating points}
\label{supp:locked}

\paragraph{Reference FROC.}
A mark matches a lesion when its squared distance is strictly below the squared half-diameter, and a missing diameter takes the default of $10\mm$ of the release. Each lesion is credited with the highest score among its matching marks, marks that match only excluded findings are ignored, and further marks on an already credited lesion are not counted as separate detections. The V2 stream has \EvalMarks{} proposals, and the reference has \EvalLesions{} lesions, of which \ProposedLesions{} are matched by at least one proposal without the mark cap. The released evaluator reproduces our sensitivities and counts for representative score streams of every comparison, both with the 100-mark cap and without it. Because the cap depends on the scores, it can change which lesion matches remain.

In each patient-level bootstrap resample, the FROC curve is recomputed from the best matching score of every lesion and all false-positive scores of the complete scan stream, including ties, and the sensitivity at each rate uses the interpolation of the release. Paired differences compare both models at the same rate within a resample. The CPM differences of band versus ball and of the two compact-head contrasts (endpoints versus scalar at fixed logit, and embedding versus logit for the endpoint gain) form the family of three comparisons specified before evaluation; for seed~1, the primary seed, their intervals are Bonferroni-adjusted 98.33\% percentile intervals. All other intervals are pointwise 95\% percentile intervals from 2,000 resamples with seed 93041. Resamples pair the models within patients and keep the trained weights fixed. The recovery fraction divides by the full-minus-ball CPM gap in every resample in which this gap is positive.

\paragraph{Cutoffs.}
The empirical cutoff is the strict threshold above the $(n+1)$-th highest false-positive score in the evaluation stream, where $n$ is the number of scans; with ties, this can leave fewer than $n$ false marks. For transferred cutoffs, the same rule is applied to the separate calibration role and the resulting threshold is used on the evaluation stream. Cutoffs are set independently for each score and cap setting, and the achieved FP burdens are reported with the sensitivities. Ranking comparisons at equal burden use the matched $S_1$. Tab.~\ref{tab:transfer} reports the transferred cutoffs for the image rankers, and the released aggregate file contains all absolute, paired and threshold results.

\paragraph{Compact-head controls and the full-field comparison.}
On the development role, conditioning on the logit was selected by the average endpoint gain over the two image seeds. The strongest grey-level morphology score on top of the full-field ranker, by average CPM, was the positive mean of the grey-level shell, ahead of the grey-level maximum and the unpaired alternatives. This shell score summarizes greyscale dilation around a bright component inside the lung found in the image~\cite{leemput2015shell}, and the conditional paired and shuffled-pair heads share its frozen base score. We regarded the paired head as practically useful on top of the full field if it gained at least one point in both CPM and matched $S_1$ with positive lower bounds, compared favourably with the pairing controls, and had no interval lower endpoint below $-2$ sensitivity points at the three lowest FP rates. For the increment of the paired head over the selected morphology score, seed~1 stays within this margin, whereas seed~2 shows possible harm at 0.125 FP/scan, and both CPM upper bounds lie below the $+1$ point threshold. Tabs.~\ref{tab:locked_factorial}--\ref{tab:locked_practical} separate this bounded result from the effects of representation and conditioning. The tables show the trained heads throughout; the development choice of a zero residual is recorded separately in the aggregate file.

\section{Follow-up comparisons and seed-2 attention heads}
\label{supp:new}

The native-source condition exposes additional native voxels to the image network; it is trained under its own visibility with the same architecture, optimizer and checkpoint rule. The direction and width comparisons in Tabs.~\ref{tab:attention_gains}--\ref{tab:attention_absolute} use exactly the same 384 interpolated samples as the compact head. The two comparisons thus answer different questions: the former changes native coverage at a fixed point of entry into the image network, and the latter isolates angular structure and capacity at fixed sparse access and score fusion. Each sparse result uses separate head and image weights per seed, whereas the single-stream and native-source models use one seed.

For seed~2, the direction-blind heads of widths 48 and 96 and the direction-aware head of width 96 selected their first evaluated checkpoint (epoch 4) and reproduce the CPM of the frozen ball-only ranker to numerical precision. All seven pooled sensitivities coincide, but their 301,306 aligned proposal scores differ from those of the ball-only ranker, as do the rankings and ties; epoch 4 is not the zero-residual initialization. The paired bootstrap uses shared patient resamples and recomputes the FROC curve of each model, so the nonzero difference intervals in Tabs.~\ref{tab:attention_contrasts}--\ref{tab:attention_absolute} are compatible with equal pooled CPM. We report these intervals as computed, and the seed-2 geometry and width contrasts are inconclusive.

For the single-stream comparison, the full-field model had to reach 80\% CPM on both development roles to serve as a competent reference; Tab.~\ref{tab:new_dev} reports both. The CPM gain of the single-stream model has a pointwise interval above zero, whereas its gain at 1~FP/scan is imprecise. The aggregate file retains, for all models, the capped and uncapped results, the results on both halves of the evaluation pool, all seven matched rates and the transferred cutoffs.

\section{Anatomy groups over all proposals}
\label{supp:contour}

We verified the series and study crosswalk, the slice positions of the reader contours and the native physical orientation. Inclusion polygons are rasterized on the native grid, and the corresponding exclusion polygons of the same reader are removed. Spatial consistency links each reader's lesion to a unique reference finding, and a voxel counts as lesion tissue if it lies in strictly more than half of the linked masks. Every linked evaluation lesion has at least two usable readers. Exact strict matches were recomputed from the release coordinates before any score was consulted, and all 604 reference--proposal pairs were reconciled with the archived contour inventory. The 598 pairs with a contour link give the same overlap labels for the native band and for the native sources that can reach the centred crop. A lesion extends into the band if its tissue overlaps the band of any of its matching proposals, regardless of their scores.

In each patient resample, the FROC axis is given by the false-positive count over all scans and the scores of the respective model. The event counts of the four groups use identical ties and interpolation abscissae, so the interpolated group counts sum to the full FROC numerator at every rate. A group contributes a sensitivity only in resamples that contain at least one of its lesions; the two small groups without a matching proposal or a contour link are descriptive, and resamples without a lesion of the respective group are omitted from its interval. The six linked lesions without a matching proposal always count as misses and remain in the full denominator. The difference between the groups that extend into and do not reach the band is a paired difference of subgroup gains, each with its own denominator.

The per-lesion decision analysis of Tab.~\ref{tab:decision} is descriptive. At a fixed cutoff, a lesion is credited to its highest-scoring retained matching mark, and this mark can differ between models. The analysis records whether the band of this decision mark intersects the contour, whereas the grouping above depends only on anatomy and not on scores. Changes among false marks indicate changes in bright support along the half-axes; without adjudicated vascular labels, these measurements describe attenuation patterns only.

\section{Earlier comparisons and reproducibility}
\label{supp:repro}

Earlier comparisons with Hessian-based readouts and visibility on this benchmark are reported for context in Tab.~\ref{tab:historical}. The Hessian procedure samples opposite half-axes through an interior anchor selected by the Hessian, with the sampling spheres still centred at the proposal, so its pairing result concerns that classical ranker and construction. We verified all seven earlier rate estimates and their intervals against the saved complete score streams, and the original 3D ball-only and full-field checkpoints and their development scores against the earlier archived fits.

The anonymized source code includes a self-contained renderer for all figures and tables and the aggregate inputs, which contain the evaluated rates, effects, uncertainty estimates, fit summaries and denominators needed to regenerate them. Private execution manifests record the exact CT and cache files, the released coordinates, the patient-disjoint roles, the selected checkpoints and the exported scores. The native-source, spatial and architecture comparisons fixed their choices before calibration and evaluation, and the anatomy groups were fixed before scores were linked to them. Model thresholds, proposal streams and raw reader records are kept separately from the review package.

\section{Exploratory size/displacement audit}
\label{supp:sizeaudit}

The audit crosses the diameter bins $<6$, 6--$<10$ and $\ge10\mm$ with the bins $<1.5$, 1.5--$<3$ and $\ge3\mm$ of the minimum displacement of a matching proposal; these bins were fixed before the outcomes were computed. It is a post hoc analysis of reused data, not a prospective preregistration. Displacement is the minimum over all matching proposals and is not selected by the winning score. All saved geometry and decision hashes and the joins of annotations to patients were verified; 449 lesions with a valid contour have matching proposals, and 72 of them extend into the band. Decisions at the calibration-fixed thresholds are primary, and the analysis at empirical thresholds is reported separately. Over all 461 lesions, the detections of seed~1 go from 397 to 414 at calibration-fixed thresholds (28 gains, 11 losses) and from 395 to 414 at empirical thresholds (28 gains, 9 losses); those of seed~2 go from 398 to 409 and from 395 to 411, respectively.

For the 449 lesions, the calibration-fixed gains are $+3.79\,[1.13,\,6.51]$ and $+2.45\,[-0.23,\,5.08]$ points for the two seeds. Lesions that extend into the band improve from 51 to 61 detections for both seeds ($+13.89\,[4.48,\,23.44]$ and $+13.89\,[5.00,\,22.73]$), and the other lesions from 340 to 347 and from 341 to 342 ($+1.86\,[-0.76,\,4.49]$ and $+0.27\,[-2.36,\,2.79]$). Lesions that extend into the band have a median of 3 matching proposals versus 1 and a minimum displacement of $1.494$ versus $0.657\mm$.

Tab.~\ref{tab:size_displacement} reports every cell, including empty and single-lesion cells. Both cells with substantial numbers of lesions in each group contain lesions of at least $10\mm$. At a displacement below $1.5\mm$, 24 lesions do not reach and 36 extend into the band (median proposal counts 1 and 3.5); at 1.5--$<3\mm$, the counts are 11 and 30 (medians 1 and 2). In the first cell, lesions extending into the band gain $13.89\,[0,\,28.21]$ and $5.56\,[-5.00,\,17.86]$ points and the other lesions gain nothing; in the second cell, they gain $13.33\,[-4.17,\,30.77]$ and $20.00\,[3.45,\,36.00]$ points, versus 0 and $9.09\,[0,\,30.77]$. These cells neither show an effect that is separated for both seeds nor remove the confounding by proposal count. Intervals use 4,000 paired percentile resamples of the 353 patients, conditional on fixed scores and cutoffs and unadjusted for multiplicity, and resamples with a zero denominator are omitted. We estimate no adjusted causal effect, because the common support is sparse and the proposal counts remain imbalanced.

\section{Further discussion and limitations}
\label{supp:discussion}

\paragraph{Coverage versus continuity.}
Sparse access recovers considerably less than the band, regardless of how it is read: exposing the exact voxels behind the 384 samples to the image network closes less than half of the gap, and attention over the sample directions adds little to a late readout. Observing a few well-chosen locations, even together with their geometry, is therefore no substitute for observing a contiguous region. Because the native-source condition changes the number and the contiguity of the visible voxels together, varying their arrangement at a fixed count is the natural next experiment to separate the two.

\paragraph{What the band adds.}
The band is defined around the proposal rather than the lesion, so large lesions and off-centre proposals place lesion tissue inside it, and these are the cases with the largest gains under a grouping fixed before scoring. A $10\mm$ ball contains most small nodules entirely: all lesions below $6\mm$ belong to the group that does not reach the band, whose gain at 1~FP/scan is small. For lesions of at least $10\mm$, the ball cuts through the nodule, and two thirds of them extend into the band. Large lesions also attract more proposals: the 72 lesions that extend into the band have 221 matching proposals, whereas every other lesion with a proposal has exactly one, and many of these proposals lie away from the lesion centre. The advantage reported for larger context may therefore partly be a coverage advantage for large lesions, which suggests stratifying analyses of context by lesion size and choosing patch sizes from the size distribution of the target lesions. Small benefits from beyond the annotated lesion remain compatible with the data, especially at higher FP rates, and part of the FROC change may come from better rejection of false proposals rather than a changed appearance of true ones; attributing gains to vessels, attachments or lesion margins would require further interventions or anatomical labels.

\paragraph{Limitations.}
Masking introduces artificial air boundaries to which each model can adapt, and the native-source and band conditions differ not only in the voxels they expose but also in how these boundaries appear after resampling. We studied a single band, chosen to enclose the support of the samples; other radii and equal-volume regions remain to be explored. The sparse heads use unoriented axis pairs and a single fusion point, so their results are specific to this design, and since three seed-2 attention heads tie the ball-only ranker in CPM (Sec.~\ref{supp:new}), the evidence on sample geometry rests mainly on seed~1. The native-source and single-stream comparisons use one training seed each, and patient-level intervals do not capture variation across training runs. Finally, all models share one candidate generator on a reused benchmark. A band condition with the contoured lesion voxels set to air would turn the anatomy analysis into a direct intervention, and an untouched full-scan cohort with a fixed proposal stream would test the generality of our findings.

\onecolumn
\section{Detailed results tables}
\label{supp:tables}
The following tables report all absolute results, paired effects and controls. Unless stated otherwise, CPM and sensitivities are in percent, effects are in points, and intervals are pointwise 95\% patient-level bootstrap intervals.

\subsection*{Absolute visibility results and paired operating points}
\begin{table*}[t]
\centering
\small
\caption{Visibility results on the 353 evaluation scans. Absolute CPM is in percent; band-minus-ball effects are in points. $^\dagger$: interval adjusted for the three comparisons specified before evaluation; all other intervals are pointwise 95\%. Recovery is the band gain divided by the full-minus-ball CPM gap.}
\label{tab:visibility_supp}
\begin{tabular}{@{}lrrrrrr@{}}
\toprule
Architecture & Ball & Ball + band & Full & $\Delta$ CPM & $\Delta S_1$ & Recovery \\
\midrule
Two-scale 3D (seed 1) & 83.58 & 87.60 & 89.37 & $+4.03\,[1.85,\,6.17]$$^\dagger$ & $+4.12\,[1.45,\,6.71]$ & 70\% \\
Two-scale 3D (seed 2) & 82.83 & 87.05 & 88.63 & $+4.21\,[2.51,\,6.17]$ & $+3.47\,[0.67,\,6.59]$ & 73\% \\
Single-stream 3D (seed 1) & 80.94 & 83.85 & 85.84 & $+2.91\,[0.69,\,5.48]$ & $+0.87\,[-2.59,\,4.73]$ & 59\% \\
Three-view 2D (seed 2) & 81.28 & 82.55 & 84.75 & $+1.27\,[-0.71,\,3.58]$ & $+0.22\,[-2.49,\,2.88]$ & 37\% \\
\bottomrule
\end{tabular}
\end{table*}

\begin{table*}[t]
\centering
\small
\caption{Sensitivity (\%) at 1~FP/scan by anatomy group, computed from one FROC curve per model over all 461 lesions and all false marks. Effects are paired band-minus-ball differences in points with pointwise 95\% intervals.}
\label{tab:anatomy_one}
\begin{tabular}{@{}llrrrr@{}}
\toprule
Seed & Group & $n$ & Ball & Band & $\Delta S_1$ \\
\midrule
1 & Does not reach band & 377 & 89.92 & 92.04 & $+2.12\,[-0.52,\,4.79]$ \\
1 & Extends into band & 72 & 69.44 & 84.72 & $+15.28\,[5.81,\,25.00]$ \\
1 & No matching proposal & 6 & 0.00 & 0.00 & $+0.00\,[0.00,\,0.00]$ \\
1 & No contour link & 6 & 100.00 & 100.00 & $+0.00\,[0.00,\,0.00]$ \\
2 & Does not reach band & 377 & 89.92 & 90.98 & $+1.06\,[-1.68,\,4.27]$ \\
2 & Extends into band & 72 & 69.44 & 86.11 & $+16.67\,[6.85,\,26.09]$ \\
2 & No matching proposal & 6 & 0.00 & 0.00 & $+0.00\,[0.00,\,0.00]$ \\
2 & No contour link & 6 & 100.00 & 100.00 & $+0.00\,[0.00,\,0.00]$ \\
\bottomrule
\end{tabular}
\end{table*}

\begin{table}[H]
\centering
\footnotesize\setlength{\tabcolsep}{4pt}
\caption{Absolute CPM and sensitivity (\%) at the seven FP/scan rates of the CPM. The native-source condition uses the two-scale architecture and seed~1.}
\label{tab:absolute}
\begin{tabular}{@{}llrrrrrrrr@{}}
\toprule
Fit & Input & CPM & 0.125 & 0.25 & 0.5 & 1 & 2 & 4 & 8 \\
\midrule
Two-scale, seed 1 & Ball & 83.58 & 70.07 & 76.14 & 81.78 & 85.68 & 88.07 & 90.67 & 92.62 \\
Two-scale, seed 1 & Band & 87.60 & 73.97 & 80.69 & 85.47 & 89.80 & 92.41 & 94.58 & 96.31 \\
Two-scale, seed 1 & Full & 89.37 & 79.18 & 82.65 & 87.64 & 90.67 & 93.71 & 95.23 & 96.53 \\
\midrule
Two-scale, seed 2 & Ball & 82.83 & 65.51 & 74.40 & 80.04 & 85.68 & 89.59 & 91.54 & 93.06 \\
Two-scale, seed 2 & Band & 87.05 & 72.67 & 79.61 & 85.25 & 89.15 & 92.62 & 94.36 & 95.66 \\
Two-scale, seed 2 & Full & 88.63 & 77.01 & 80.69 & 87.20 & 90.46 & 93.93 & 95.01 & 96.10 \\
\midrule
Three-view, seed 2 & Ball & 81.28 & 66.59 & 72.02 & 77.87 & 82.86 & 86.98 & 90.24 & 92.41 \\
Three-view, seed 2 & Band & 82.55 & 68.98 & 74.40 & 79.61 & 83.08 & 87.85 & 90.89 & 93.06 \\
Three-view, seed 2 & Full & 84.75 & 69.20 & 75.92 & 82.00 & 86.33 & 91.11 & 93.28 & 95.44 \\
\midrule
Single-stream, seed 1 & Ball & 80.94 & 65.08 & 70.28 & 77.87 & 83.95 & 88.07 & 89.80 & 91.54 \\
Single-stream, seed 1 & Band & 83.85 & 70.50 & 75.49 & 79.39 & 84.82 & 90.02 & 92.19 & 94.58 \\
Single-stream, seed 1 & Full & 85.84 & 71.58 & 79.18 & 83.73 & 87.64 & 90.67 & 92.84 & 95.23 \\
\midrule
Native sources, seed 1 & Native sources & 85.31 & 70.72 & 77.22 & 83.51 & 88.29 & 90.46 & 92.62 & 94.36 \\
\bottomrule
\end{tabular}
\end{table}

\begin{table}[H]
\centering
\footnotesize\setlength{\tabcolsep}{4pt}
\caption{Paired sensitivity differences (points) at each FP/scan rate. Each cell gives the estimate above its pointwise 95\% patient-level interval.}
\label{tab:paired_rates}
\begin{tabular}{@{}lrrrrrrr@{}}
\toprule
Contrast & 0.125 & 0.25 & 0.5 & 1 & 2 & 4 & 8 \\
\midrule
Two-scale, seed 1 / Band $-$ ball & \raisebox{-.4\height}{\shortstack{$+3.9$\\$[0.0,\,8.4]$}} & \raisebox{-.4\height}{\shortstack{$+4.6$\\$[-0.2,\,7.9]$}} & \raisebox{-.4\height}{\shortstack{$+3.7$\\$[0.0,\,6.7]$}} & \raisebox{-.4\height}{\shortstack{$+4.1$\\$[1.4,\,6.7]$}} & \raisebox{-.4\height}{\shortstack{$+4.3$\\$[2.4,\,6.7]$}} & \raisebox{-.4\height}{\shortstack{$+3.9$\\$[2.1,\,6.1]$}} & \raisebox{-.4\height}{\shortstack{$+3.7$\\$[1.7,\,5.8]$}} \\
Two-scale, seed 1 / Full $-$ ball & \raisebox{-.4\height}{\shortstack{$+9.1$\\$[4.8,\,13.0]$}} & \raisebox{-.4\height}{\shortstack{$+6.5$\\$[2.7,\,10.3]$}} & \raisebox{-.4\height}{\shortstack{$+5.9$\\$[3.2,\,8.5]$}} & \raisebox{-.4\height}{\shortstack{$+5.0$\\$[2.6,\,7.7]$}} & \raisebox{-.4\height}{\shortstack{$+5.6$\\$[3.1,\,7.9]$}} & \raisebox{-.4\height}{\shortstack{$+4.6$\\$[2.5,\,6.8]$}} & \raisebox{-.4\height}{\shortstack{$+3.9$\\$[2.1,\,6.0]$}} \\
Two-scale, seed 1 / Full $-$ band & \raisebox{-.4\height}{\shortstack{$+5.2$\\$[0.4,\,8.7]$}} & \raisebox{-.4\height}{\shortstack{$+2.0$\\$[-0.2,\,5.8]$}} & \raisebox{-.4\height}{\shortstack{$+2.2$\\$[-0.2,\,5.0]$}} & \raisebox{-.4\height}{\shortstack{$+0.9$\\$[-1.0,\,3.2]$}} & \raisebox{-.4\height}{\shortstack{$+1.3$\\$[-1.0,\,3.0]$}} & \raisebox{-.4\height}{\shortstack{$+0.7$\\$[-1.3,\,2.5]$}} & \raisebox{-.4\height}{\shortstack{$+0.2$\\$[-0.9,\,1.6]$}} \\
\midrule
Two-scale, seed 2 / Band $-$ ball & \raisebox{-.4\height}{\shortstack{$+7.2$\\$[1.7,\,13.0]$}} & \raisebox{-.4\height}{\shortstack{$+5.2$\\$[2.3,\,9.3]$}} & \raisebox{-.4\height}{\shortstack{$+5.2$\\$[2.4,\,8.6]$}} & \raisebox{-.4\height}{\shortstack{$+3.5$\\$[0.7,\,6.6]$}} & \raisebox{-.4\height}{\shortstack{$+3.0$\\$[0.9,\,5.1]$}} & \raisebox{-.4\height}{\shortstack{$+2.8$\\$[0.7,\,5.3]$}} & \raisebox{-.4\height}{\shortstack{$+2.6$\\$[0.7,\,4.4]$}} \\
Two-scale, seed 2 / Full $-$ ball & \raisebox{-.4\height}{\shortstack{$+11.5$\\$[6.7,\,18.4]$}} & \raisebox{-.4\height}{\shortstack{$+6.3$\\$[3.3,\,11.3]$}} & \raisebox{-.4\height}{\shortstack{$+7.2$\\$[3.8,\,10.6]$}} & \raisebox{-.4\height}{\shortstack{$+4.8$\\$[2.4,\,8.0]$}} & \raisebox{-.4\height}{\shortstack{$+4.3$\\$[2.1,\,6.6]$}} & \raisebox{-.4\height}{\shortstack{$+3.5$\\$[1.4,\,5.9]$}} & \raisebox{-.4\height}{\shortstack{$+3.0$\\$[1.1,\,5.0]$}} \\
Two-scale, seed 2 / Full $-$ band & \raisebox{-.4\height}{\shortstack{$+4.3$\\$[0.0,\,10.1]$}} & \raisebox{-.4\height}{\shortstack{$+1.1$\\$[-1.8,\,5.3]$}} & \raisebox{-.4\height}{\shortstack{$+2.0$\\$[-1.0,\,4.4]$}} & \raisebox{-.4\height}{\shortstack{$+1.3$\\$[-0.6,\,3.7]$}} & \raisebox{-.4\height}{\shortstack{$+1.3$\\$[-0.2,\,2.8]$}} & \raisebox{-.4\height}{\shortstack{$+0.7$\\$[-0.7,\,2.2]$}} & \raisebox{-.4\height}{\shortstack{$+0.4$\\$[-0.5,\,1.6]$}} \\
\midrule
Three-view, seed 2 / Band $-$ ball & \raisebox{-.4\height}{\shortstack{$+2.4$\\$[-2.9,\,8.6]$}} & \raisebox{-.4\height}{\shortstack{$+2.4$\\$[-1.9,\,7.1]$}} & \raisebox{-.4\height}{\shortstack{$+1.7$\\$[-1.6,\,4.8]$}} & \raisebox{-.4\height}{\shortstack{$+0.2$\\$[-2.5,\,2.9]$}} & \raisebox{-.4\height}{\shortstack{$+0.9$\\$[-2.1,\,3.3]$}} & \raisebox{-.4\height}{\shortstack{$+0.7$\\$[-1.2,\,2.8]$}} & \raisebox{-.4\height}{\shortstack{$+0.7$\\$[-0.9,\,2.8]$}} \\
Three-view, seed 2 / Full $-$ ball & \raisebox{-.4\height}{\shortstack{$+2.6$\\$[-2.9,\,8.8]$}} & \raisebox{-.4\height}{\shortstack{$+3.9$\\$[0.0,\,9.3]$}} & \raisebox{-.4\height}{\shortstack{$+4.1$\\$[0.9,\,7.9]$}} & \raisebox{-.4\height}{\shortstack{$+3.5$\\$[0.9,\,6.6]$}} & \raisebox{-.4\height}{\shortstack{$+4.1$\\$[1.6,\,6.5]$}} & \raisebox{-.4\height}{\shortstack{$+3.0$\\$[1.1,\,5.4]$}} & \raisebox{-.4\height}{\shortstack{$+3.0$\\$[1.3,\,5.2]$}} \\
Three-view, seed 2 / Full $-$ band & \raisebox{-.4\height}{\shortstack{$+0.2$\\$[-5.8,\,5.4]$}} & \raisebox{-.4\height}{\shortstack{$+1.5$\\$[-1.7,\,5.7]$}} & \raisebox{-.4\height}{\shortstack{$+2.4$\\$[0.0,\,5.7]$}} & \raisebox{-.4\height}{\shortstack{$+3.3$\\$[1.0,\,6.2]$}} & \raisebox{-.4\height}{\shortstack{$+3.3$\\$[1.1,\,5.9]$}} & \raisebox{-.4\height}{\shortstack{$+2.4$\\$[0.6,\,4.2]$}} & \raisebox{-.4\height}{\shortstack{$+2.4$\\$[0.4,\,4.2]$}} \\
\midrule
Single-stream, seed 1 / Band $-$ ball & \raisebox{-.4\height}{\shortstack{$+5.4$\\$[0.6,\,12.8]$}} & \raisebox{-.4\height}{\shortstack{$+5.2$\\$[1.1,\,9.3]$}} & \raisebox{-.4\height}{\shortstack{$+1.5$\\$[-2.6,\,5.3]$}} & \raisebox{-.4\height}{\shortstack{$+0.9$\\$[-2.6,\,4.7]$}} & \raisebox{-.4\height}{\shortstack{$+2.0$\\$[-1.0,\,4.5]$}} & \raisebox{-.4\height}{\shortstack{$+2.4$\\$[0.0,\,4.6]$}} & \raisebox{-.4\height}{\shortstack{$+3.0$\\$[1.1,\,5.2]$}} \\
Single-stream, seed 1 / Full $-$ ball & \raisebox{-.4\height}{\shortstack{$+6.5$\\$[2.6,\,14.0]$}} & \raisebox{-.4\height}{\shortstack{$+8.9$\\$[4.2,\,13.4]$}} & \raisebox{-.4\height}{\shortstack{$+5.9$\\$[2.2,\,10.2]$}} & \raisebox{-.4\height}{\shortstack{$+3.7$\\$[0.4,\,7.0]$}} & \raisebox{-.4\height}{\shortstack{$+2.6$\\$[0.2,\,5.2]$}} & \raisebox{-.4\height}{\shortstack{$+3.0$\\$[0.7,\,5.3]$}} & \raisebox{-.4\height}{\shortstack{$+3.7$\\$[1.6,\,6.0]$}} \\
Single-stream, seed 1 / Full $-$ band & \raisebox{-.4\height}{\shortstack{$+1.1$\\$[-3.0,\,6.5]$}} & \raisebox{-.4\height}{\shortstack{$+3.7$\\$[-0.2,\,7.2]$}} & \raisebox{-.4\height}{\shortstack{$+4.3$\\$[1.7,\,8.1]$}} & \raisebox{-.4\height}{\shortstack{$+2.8$\\$[-0.7,\,5.8]$}} & \raisebox{-.4\height}{\shortstack{$+0.7$\\$[-1.4,\,3.2]$}} & \raisebox{-.4\height}{\shortstack{$+0.7$\\$[-1.4,\,2.9]$}} & \raisebox{-.4\height}{\shortstack{$+0.7$\\$[-1.1,\,2.4]$}} \\
\midrule
Native sources, seed 1 / Sources $-$ ball & \raisebox{-.4\height}{\shortstack{$+0.7$\\$[-6.0,\,4.3]$}} & \raisebox{-.4\height}{\shortstack{$+1.1$\\$[-2.8,\,4.6]$}} & \raisebox{-.4\height}{\shortstack{$+1.7$\\$[-1.3,\,4.1]$}} & \raisebox{-.4\height}{\shortstack{$+2.6$\\$[-0.6,\,4.6]$}} & \raisebox{-.4\height}{\shortstack{$+2.4$\\$[0.0,\,4.3]$}} & \raisebox{-.4\height}{\shortstack{$+2.0$\\$[0.2,\,3.8]$}} & \raisebox{-.4\height}{\shortstack{$+1.7$\\$[0.2,\,3.4]$}} \\
Native sources, seed 1 / Band $-$ sources & \raisebox{-.4\height}{\shortstack{$+3.3$\\$[-0.7,\,10.4]$}} & \raisebox{-.4\height}{\shortstack{$+3.5$\\$[-0.8,\,6.5]$}} & \raisebox{-.4\height}{\shortstack{$+2.0$\\$[-0.9,\,5.1]$}} & \raisebox{-.4\height}{\shortstack{$+1.5$\\$[-0.5,\,4.7]$}} & \raisebox{-.4\height}{\shortstack{$+2.0$\\$[0.4,\,4.4]$}} & \raisebox{-.4\height}{\shortstack{$+2.0$\\$[0.4,\,3.7]$}} & \raisebox{-.4\height}{\shortstack{$+2.0$\\$[0.2,\,3.6]$}} \\
\bottomrule
\end{tabular}
\end{table}


\clearpage
\subsection*{Consistency checks and transferred cutoffs}
\begin{table}[H]
\centering
\small
\caption{Consistency checks: band minus ball for the dense visibility conditions (native sources minus ball for the source condition) on each half of the evaluation pool and without the 100-mark cap. Intervals are pointwise 95\%.}
\label{tab:consistency}
\begin{tabular}{@{}llrrr@{}}
\toprule
Fit & Scope & Scans & $\Delta$ CPM & $\Delta S_1$ \\
\midrule
Two-scale, seed 1 & Subsets 6--7 & 177 & $+4.38\,[1.96,\,7.45]$ & $+4.62\,[1.89,\,7.66]$ \\
Two-scale, seed 1 & Pooled & 353 & $+4.03\,[2.16,\,5.71]$ & $+4.12\,[1.45,\,6.71]$ \\
Two-scale, seed 1 & Pooled, uncapped & 353 & $+4.03\,[2.16,\,5.71]$ & $+4.12\,[1.45,\,6.71]$ \\
Two-scale, seed 1 & Subsets 8--9 & 176 & $+3.65\,[0.84,\,6.64]$ & $+3.14\,[-0.90,\,7.69]$ \\
\midrule
Two-scale, seed 2 & Subsets 6--7 & 177 & $+6.18\,[4.01,\,8.73]$ & $+6.72\,[3.32,\,10.96]$ \\
Two-scale, seed 2 & Pooled & 353 & $+4.21\,[2.51,\,6.17]$ & $+3.47\,[0.67,\,6.59]$ \\
Two-scale, seed 2 & Pooled, uncapped & 353 & $+4.21\,[2.52,\,6.19]$ & $+3.47\,[0.67,\,6.59]$ \\
Two-scale, seed 2 & Subsets 8--9 & 176 & $+3.14\,[0.38,\,5.64]$ & $+0.90\,[-3.79,\,4.04]$ \\
\midrule
Three-view, seed 2 & Pooled & 353 & $+1.27\,[-0.71,\,3.58]$ & $+0.22\,[-2.49,\,2.88]$ \\
Three-view, seed 2 & Pooled, uncapped & 353 & $+1.27\,[-0.71,\,3.58]$ & $+0.22\,[-2.49,\,2.88]$ \\
\midrule
Single-stream, seed 1 & Subsets 6--7 & 177 & $+3.00\,[0.25,\,6.51]$ & $+3.36\,[-1.71,\,7.02]$ \\
Single-stream, seed 1 & Pooled & 353 & $+2.91\,[0.69,\,5.48]$ & $+0.87\,[-2.59,\,4.73]$ \\
Single-stream, seed 1 & Pooled, uncapped & 353 & $+2.91\,[0.69,\,5.48]$ & $+0.87\,[-2.59,\,4.73]$ \\
Single-stream, seed 1 & Subsets 8--9 & 176 & $+2.37\,[-0.97,\,6.05]$ & $+0.45\,[-5.71,\,6.05]$ \\
\midrule
Native sources, seed 1 & Subsets 6--7 & 177 & $+1.62\,[-1.05,\,3.72]$ & $+2.94\,[-1.02,\,5.85]$ \\
Native sources, seed 1 & Pooled & 353 & $+1.74\,[-0.21,\,3.12]$ & $+2.60\,[-0.63,\,4.64]$ \\
Native sources, seed 1 & Pooled, uncapped & 353 & $+1.74\,[-0.21,\,3.12]$ & $+2.60\,[-0.63,\,4.64]$ \\
Native sources, seed 1 & Subsets 8--9 & 176 & $+1.92\,[-0.59,\,4.40]$ & $+1.79\,[-2.15,\,5.63]$ \\
\bottomrule
\end{tabular}
\end{table}

\begin{table*}[!t]
\centering
\small
\caption{Thresholds transferred from the separate 88-scan calibration role. Entries give ball/band (or ball/native sources); sensitivity is in percent, differences in points, FP burden in marks per scan. Intervals are pointwise 95\%.}
\label{tab:transfer}
\begin{tabular}{@{}lrrrr@{}}
\toprule
Fit & Sensitivity & Achieved FP/scan & $\Delta$ Sens. & $\Delta$ FP/scan \\
\midrule
Two-scale, seed 1 & 86.12 / 89.80 & 1.091 / 1.045 & $+3.69\,[1.02,\,6.53]$ & $-0.045\,[-0.187,\,0.091]$ \\
Two-scale, seed 2 & 86.33 / 88.72 & 1.093 / 0.977 & $+2.39\,[-0.22,\,5.03]$ & $-0.116\,[-0.238,\,-0.003]$ \\
Three-view, seed 2 & 82.43 / 83.73 & 0.960 / 1.195 & $+1.30\,[-1.32,\,3.86]$ & $+0.235\,[0.105,\,0.363]$ \\
Single-stream, seed 1 & 82.65 / 83.95 & 0.898 / 0.901 & $+1.30\,[-2.14,\,4.72]$ & $+0.003\,[-0.127,\,0.133]$ \\
Native sources, seed 1 & 86.12 / 86.77 & 1.091 / 0.924 & $+0.65\,[-1.73,\,3.16]$ & $-0.167\,[-0.306,\,-0.025]$ \\
\bottomrule
\end{tabular}
\end{table*}


\clearpage
\subsection*{Compact-head factorial, pairing controls and full-field comparisons}
\begin{table*}[!t]
\centering
\footnotesize\setlength{\tabcolsep}{4pt}
\caption{Compact-head factorial (initial fits), each pathway against its own image-only baseline. Effects are in points. $^\dagger$: 98.33\% interval of the three-comparison family specified before evaluation; all other intervals are pointwise 95\%.}
\label{tab:locked_factorial}
\begin{tabular}{@{}lrrrr@{}}
\toprule
Contrast & Seed 1 $\Delta$ CPM & Seed 1 $\Delta S_1$ & Seed 2 $\Delta$ CPM & Seed 2 $\Delta S_1$ \\
\midrule
Embedding-only readout $-$ ball & $+0.00\,[-0.10,\,0.09]$ & $+0.00\,[-0.23,\,0.00]$ & $-0.03\,[-0.10,\,0.09]$ & $+0.00\,[0.00,\,0.00]$ \\
Scalar gain, logit & $+0.25\,[-0.37,\,0.47]$ & $+0.00\,[-1.17,\,0.88]$ & $+0.31\,[-0.00,\,1.04]$ & $+0.00\,[-0.47,\,1.30]$ \\
Scalar gain, embedding & $+0.28\,[-0.38,\,0.59]$ & $+0.65\,[-1.04,\,1.72]$ & $+0.34\,[-0.06,\,1.00]$ & $+0.00\,[-0.63,\,1.64]$ \\
Endpoint gain, logit & $+0.87\,[0.16,\,1.30]$ & $+1.52\,[-0.22,\,2.68]$ & $+0.12\,[-0.25,\,1.05]$ & $+0.22\,[-0.76,\,1.80]$ \\
Endpoint gain, embedding & $+0.77\,[-0.26,\,1.60]$ & $+0.65\,[-0.89,\,2.34]$ & $+0.06\,[-1.09,\,1.43]$ & $+0.65\,[-1.70,\,3.08]$ \\
Endpoints $-$ scalar, logit & $+0.62\,[-0.09,\,1.38]$$^\dagger$ & $+1.52\,[-0.21,\,2.99]$ & $-0.19\,[-0.72,\,0.60]$ & $+0.22\,[-1.08,\,1.44]$ \\
Endpoints $-$ scalar, embedding & $+0.50\,[-0.41,\,1.55]$ & $+0.00\,[-1.18,\,2.09]$ & $-0.28\,[-1.53,\,0.97]$ & $+0.65\,[-1.96,\,2.75]$ \\
Embedding $-$ logit, scalar & $+0.03\,[-0.33,\,0.42]$ & $+0.65\,[-0.64,\,1.75]$ & $+0.03\,[-0.33,\,0.30]$ & $+0.00\,[-0.87,\,1.00]$ \\
Embedding $-$ logit, endpoints & $-0.09\,[-1.23,\,1.05]$$^\dagger$ & $-0.87\,[-2.29,\,1.16]$ & $-0.06\,[-1.39,\,0.91]$ & $+0.43\,[-1.95,\,2.46]$ \\
Representation $\times$ input interaction & $-0.12\,[-1.13,\,0.89]$ & $-1.52\,[-3.23,\,1.08]$ & $-0.09\,[-1.37,\,0.93]$ & $+0.43\,[-1.95,\,2.41]$ \\
Band $-$ selected sparse gain & $+3.16\,[1.46,\,5.04]$ & $+2.60\,[0.41,\,5.70]$ & $+4.09\,[2.12,\,5.78]$ & $+3.25\,[0.19,\,6.13]$ \\
\bottomrule
\end{tabular}
\end{table*}

\begin{table*}[!t]
\centering
\footnotesize\setlength{\tabcolsep}{4pt}
\caption{Paired compact head (initial fit) against shuffled pairs, unpaired half-axes and grey-level morphology heads. Shuffled-pair values average eight evaluation permutations of one fitted head. Effects in points with pointwise 95\% intervals.}
\label{tab:locked_compact}
\begin{tabular}{@{}lrrrr@{}}
\toprule
Contrast & Seed 1 $\Delta$ CPM & Seed 1 $\Delta S_1$ & Seed 2 $\Delta$ CPM & Seed 2 $\Delta S_1$ \\
\midrule
Paired $-$ shuffled pairs (logit) & $-0.03\,[-0.68,\,0.51]$ & $+0.76\,[-0.64,\,1.83]$ & $-0.11\,[-0.71,\,0.54]$ & $+0.35\,[-1.12,\,1.51]$ \\
Paired $-$ unpaired (logit) & $-0.31\,[-0.94,\,0.30]$ & $+0.87\,[-0.80,\,2.04]$ & $-0.31\,[-1.02,\,0.41]$ & $+0.22\,[-1.97,\,1.46]$ \\
Paired $-$ grey max. (logit) & $+0.81\,[0.19,\,1.33]$ & $+1.52\,[-0.22,\,2.68]$ & $-0.03\,[-0.32,\,0.98]$ & $+0.22\,[-0.70,\,1.72]$ \\
Paired $-$ grey mean (logit) & $+0.77\,[0.16,\,1.30]$ & $+1.52\,[-0.22,\,2.68]$ & $+0.03\,[-0.33,\,0.99]$ & $+0.22\,[-0.70,\,1.78]$ \\
Paired $-$ shuffled pairs (embedding) & $-0.69\,[-1.33,\,-0.11]$ & $-0.92\,[-2.00,\,0.43]$ & $-0.63\,[-1.48,\,0.14]$ & $-0.27\,[-1.97,\,0.80]$ \\
Paired $-$ unpaired (embedding) & $-0.90\,[-1.54,\,-0.09]$ & $-0.65\,[-1.78,\,0.74]$ & $-0.56\,[-1.64,\,0.34]$ & $+0.43\,[-2.00,\,1.30]$ \\
Paired $-$ grey max. (embedding) & $+0.68\,[-0.39,\,1.54]$ & $+0.43\,[-0.88,\,2.28]$ & $-0.12\,[-1.27,\,1.22]$ & $+1.08\,[-1.88,\,3.16]$ \\
Paired $-$ grey mean (embedding) & $+0.56\,[-0.38,\,1.47]$ & $+0.87\,[-0.44,\,2.49]$ & $-0.09\,[-1.30,\,1.22]$ & $+0.87\,[-2.02,\,3.04]$ \\
\bottomrule
\end{tabular}
\end{table*}

\begin{table}[H]
\centering
\footnotesize\setlength{\tabcolsep}{4pt}
\caption{Comparisons on top of the full-field ranker. The grey-level mean is the morphology score selected on the development role, and the conditional shuffled-pair head uses the same base. For the paired-head increment, the upper CPM bound is below $+1$ point for both seeds. Effects in points with pointwise 95\% intervals.}
\label{tab:locked_practical}
\begin{tabular}{@{}lrrrr@{}}
\toprule
Contrast & Seed 1 $\Delta$ CPM & Seed 1 $\Delta S_1$ & Seed 2 $\Delta$ CPM & Seed 2 $\Delta S_1$ \\
\midrule
Paired gain over full field & $+0.00\,[-0.52,\,0.42]$ & $+0.22\,[-0.85,\,1.06]$ & $+0.06\,[-0.52,\,0.56]$ & $+0.22\,[-0.23,\,1.00]$ \\
Morphology gain over full field & $+0.09\,[-0.20,\,0.19]$ & $+0.00\,[-0.22,\,0.00]$ & $+0.12\,[-0.18,\,0.34]$ & $+0.00\,[-0.40,\,0.73]$ \\
Paired $-$ selected morphology & $-0.09\,[-0.52,\,0.40]$ & $+0.22\,[-0.84,\,1.07]$ & $-0.06\,[-0.59,\,0.48]$ & $+0.22\,[-0.49,\,0.96]$ \\
Paired $-$ unpaired & $+0.03\,[-0.53,\,0.32]$ & $+0.22\,[-0.82,\,1.09]$ & $-0.22\,[-0.56,\,0.35]$ & $+0.22\,[0.00,\,0.96]$ \\
Paired $-$ shuffled pairs & $-0.05\,[-0.50,\,0.25]$ & $-0.14\,[-0.88,\,0.49]$ & $-0.17\,[-0.50,\,0.32]$ & $+0.05\,[-0.22,\,0.89]$ \\
Paired increment over morphology & $+0.22\,[-0.29,\,0.50]$ & $+0.65\,[-0.67,\,1.14]$ & $-0.19\,[-0.86,\,0.41]$ & $-0.22\,[-1.25,\,1.10]$ \\
Conditional paired $-$ shuffled & $+0.11\,[-0.41,\,0.32]$ & $+0.43\,[-0.81,\,0.70]$ & $-0.34\,[-0.94,\,0.20]$ & $-0.49\,[-1.41,\,0.74]$ \\
\bottomrule
\end{tabular}
\end{table}

\begin{table}[H]
\centering
\small
\caption{Variation over the eight shuffled-pair permutations of each fitted head. SD in points, CPM range in percent.}
\label{tab:sham_spread}
\begin{tabular}{@{}llrrr@{}}
\toprule
Seed & Shuffled-pair head & CPM SD & $S_1$ SD & CPM range \\
\midrule
1 & Ball logit & 0.04 & 0.23 & 84.44--84.54 \\
1 & Ball embedding & 0.07 & 0.30 & 84.94--85.16 \\
1 & Full logit & 0.05 & 0.11 & 89.34--89.49 \\
1 & Full + morphology & 0.05 & 0.20 & 89.53--89.65 \\
2 & Ball logit & 0.16 & 0.43 & 82.86--83.39 \\
2 & Ball embedding & 0.15 & 0.19 & 83.30--83.70 \\
2 & Full logit & 0.03 & 0.10 & 88.81--88.91 \\
2 & Full + morphology & 0.04 & 0.15 & 88.84--88.97 \\
\bottomrule
\end{tabular}
\end{table}


\clearpage
\subsection*{Follow-up development results and sparse-head comparisons}
\begin{table*}[!t]
\centering
\footnotesize\setlength{\tabcolsep}{4pt}
\caption{Development-role results of the follow-up comparisons. Each cell gives CPM / $S_1$ in percent at the selected checkpoint and on the readout-development role.}
\label{tab:new_dev}
\begin{tabular}{@{}llrrrr@{}}
\toprule
Family & Condition & Seed 1 checkpoint & Seed 1 readout dev. & Seed 2 checkpoint & Seed 2 readout dev. \\
\midrule
Single-stream & Ball & 81.79 / 80.39 & 77.69 / 82.46 & -- & -- \\
Single-stream & Band & 84.03 / 84.31 & 80.20 / 84.21 & -- & -- \\
Single-stream & Full & 84.59 / 82.35 & 81.70 / 78.95 & -- & -- \\
\midrule
Native sources & Ball & 80.95 / 82.35 & 80.95 / 84.21 & -- & -- \\
Native sources & Band & 84.87 / 82.35 & 83.21 / 85.96 & -- & -- \\
Native sources & Full & 85.15 / 84.31 & 86.47 / 89.47 & -- & -- \\
Native sources & Native sources & 82.91 / 86.27 & 82.21 / 85.96 & -- & -- \\
\midrule
Sparse & Ball & 80.95 / 82.35 & 80.95 / 84.21 & 80.67 / 84.31 & 80.95 / 84.21 \\
Sparse & Band & 84.87 / 82.35 & 83.21 / 85.96 & 87.39 / 92.16 & 85.46 / 82.46 \\
Sparse & Full & 85.15 / 84.31 & 86.47 / 89.47 & 85.15 / 90.20 & 84.96 / 87.72 \\
Sparse & Compact (initial) & 80.95 / 82.35 & 81.45 / 84.21 & 80.39 / 84.31 & 81.70 / 84.21 \\
Sparse & Compact refit & 81.23 / 82.35 & 82.46 / 84.21 & 80.67 / 84.31 & 80.95 / 84.21 \\
Sparse & Direction-blind 48 & 81.23 / 82.35 & 81.70 / 84.21 & 80.67 / 84.31 & 80.95 / 84.21 \\
Sparse & Direction-aware 48 & 80.95 / 82.35 & 82.71 / 82.46 & 81.51 / 84.31 & 80.95 / 84.21 \\
Sparse & Direction-blind 96 & 81.23 / 82.35 & 81.95 / 84.21 & 80.67 / 84.31 & 80.95 / 84.21 \\
Sparse & Direction-aware 96 & 80.95 / 82.35 & 81.70 / 82.46 & 80.67 / 84.31 & 80.95 / 84.21 \\
\bottomrule
\end{tabular}
\end{table*}

\begin{table}[H]
\centering
\footnotesize\setlength{\tabcolsep}{4pt}
\caption{Direction and width contrasts of the attention heads and comparisons with the compact head. Effects in points with pointwise 95\% intervals. Seed-2 contrasts involving the three heads that tie the ball-only ranker (Sec.~\ref{supp:new}) are inconclusive.}
\label{tab:attention_contrasts}
\begin{tabular}{@{}lrrrr@{}}
\toprule
Contrast & Seed 1 $\Delta$ CPM & Seed 1 $\Delta S_1$ & Seed 2 $\Delta$ CPM & Seed 2 $\Delta S_1$ \\
\midrule
Aware $-$ blind, width 48 & $+0.19\,[-0.56,\,0.83]$ & $-0.65\,[-2.52,\,1.04]$ & $-0.22\,[-0.46,\,0.25]$ & $+0.22\,[-0.64,\,1.10]$ \\
Aware $-$ blind, width 96 & $-0.15\,[-0.63,\,0.22]$ & $-0.87\,[-1.78,\,0.00]$ & $+0.00\,[0.00,\,0.07]$ & $+0.00\,[0.00,\,0.00]$ \\
Width 96 $-$ 48, direction-blind & $+0.15\,[-0.49,\,0.82]$ & $+0.43\,[-1.03,\,1.47]$ & $+0.00\,[-0.06,\,0.03]$ & $+0.00\,[0.00,\,0.00]$ \\
Width 96 $-$ 48, direction-aware & $-0.19\,[-0.79,\,0.44]$ & $+0.22\,[-1.65,\,1.96]$ & $+0.22\,[-0.24,\,0.46]$ & $-0.22\,[-1.10,\,0.64]$ \\
Compact refit $-$ initial & $+0.62\,[-0.15,\,1.02]$ & $+0.87\,[-0.43,\,1.89]$ & $-0.03\,[-0.80,\,0.46]$ & $-0.43\,[-1.66,\,0.89]$ \\
Direction-blind 48 $-$ compact refit & $-0.68\,[-1.10,\,0.15]$ & $-1.08\,[-1.91,\,0.69]$ & $-0.09\,[-0.61,\,0.06]$ & $+0.22\,[-0.91,\,0.68]$ \\
Direction-blind 48 $-$ compact initial & $-0.06\,[-0.76,\,0.64]$ & $-0.22\,[-1.21,\,1.53]$ & $-0.12\,[-1.05,\,0.25]$ & $-0.22\,[-1.80,\,0.76]$ \\
Direction-aware 48 $-$ compact refit & $-0.50\,[-0.95,\,0.19]$ & $-1.74\,[-2.81,\,0.00]$ & $-0.31\,[-0.80,\,0.03]$ & $+0.43\,[-0.96,\,1.00]$ \\
Direction-aware 48 $-$ compact initial & $+0.12\,[-0.62,\,0.66]$ & $-0.87\,[-2.32,\,1.07]$ & $-0.34\,[-1.23,\,0.15]$ & $+0.00\,[-1.69,\,1.09]$ \\
Direction-blind 96 $-$ compact refit & $-0.53\,[-0.88,\,0.23]$ & $-0.65\,[-1.53,\,0.89]$ & $-0.09\,[-0.61,\,0.06]$ & $+0.22\,[-0.91,\,0.68]$ \\
Direction-blind 96 $-$ compact initial & $+0.09\,[-0.44,\,0.59]$ & $+0.22\,[-0.46,\,1.57]$ & $-0.12\,[-1.05,\,0.25]$ & $-0.22\,[-1.80,\,0.76]$ \\
Direction-aware 96 $-$ compact refit & $-0.68\,[-1.18,\,0.06]$ & $-1.52\,[-2.63,\,0.38]$ & $-0.09\,[-0.61,\,0.07]$ & $+0.22\,[-0.91,\,0.68]$ \\
Direction-aware 96 $-$ compact initial & $-0.06\,[-0.68,\,0.43]$ & $-0.65\,[-1.60,\,0.88]$ & $-0.12\,[-1.04,\,0.25]$ & $-0.22\,[-1.79,\,0.76]$ \\
\bottomrule
\end{tabular}
\end{table}

\begin{table*}[!t]
\centering
\footnotesize\setlength{\tabcolsep}{4pt}
\caption{Gain of each sparse head over its ball-only ranker and remaining gap to the dense band. Paired effects in points with pointwise 95\% intervals.}
\label{tab:attention_gains}
\begin{tabular}{@{}llrrrr@{}}
\toprule
Seed & Sparse head & Head $-$ ball CPM & Head $-$ ball $S_1$ & Band $-$ head CPM & Band $-$ head $S_1$ \\
\midrule
1 & Compact (initial) & $+0.87\,[0.16,\,1.30]$ & $+1.52\,[-0.22,\,2.68]$ & $+3.16\,[1.46,\,5.04]$ & $+2.60\,[0.41,\,5.70]$ \\
1 & Compact refit & $+1.49\,[0.46,\,1.82]$ & $+2.39\,[0.43,\,3.56]$ & $+2.54\,[1.05,\,4.57]$ & $+1.74\,[-0.42,\,4.96]$ \\
1 & Direction-blind 48 & $+0.81\,[-0.03,\,1.35]$ & $+1.30\,[0.00,\,2.78]$ & $+3.22\,[1.57,\,5.06]$ & $+2.82\,[0.39,\,5.39]$ \\
1 & Direction-aware 48 & $+0.99\,[0.03,\,1.47]$ & $+0.65\,[-1.19,\,2.32]$ & $+3.04\,[1.44,\,5.04]$ & $+3.47\,[0.72,\,6.53]$ \\
1 & Direction-blind 96 & $+0.96\,[0.37,\,1.23]$ & $+1.74\,[0.43,\,2.80]$ & $+3.07\,[1.50,\,4.90]$ & $+2.39\,[0.00,\,5.25]$ \\
1 & Direction-aware 96 & $+0.81\,[0.09,\,1.11]$ & $+0.87\,[0.00,\,1.79]$ & $+3.22\,[1.58,\,5.12]$ & $+3.25\,[0.83,\,5.90]$ \\
\midrule
2 & Compact (initial) & $+0.12\,[-0.25,\,1.05]$ & $+0.22\,[-0.76,\,1.80]$ & $+4.09\,[2.12,\,5.78]$ & $+3.25\,[0.19,\,6.13]$ \\
2 & Compact refit & $+0.09\,[-0.06,\,0.61]$ & $-0.22\,[-0.68,\,0.91]$ & $+4.12\,[2.33,\,5.95]$ & $+3.69\,[0.46,\,6.47]$ \\
2 & Direction-blind 48 & $+0.00\,[-0.03,\,0.00]$ & $+0.00\,[0.00,\,0.00]$ & $+4.21\,[2.50,\,6.17]$ & $+3.47\,[0.67,\,6.59]$ \\
2 & Direction-aware 48 & $-0.22\,[-0.46,\,0.24]$ & $+0.22\,[-0.64,\,1.10]$ & $+4.43\,[2.63,\,6.39]$ & $+3.25\,[0.42,\,6.42]$ \\
2 & Direction-blind 96 & $+0.00\,[-0.06,\,0.00]$ & $+0.00\,[0.00,\,0.00]$ & $+4.21\,[2.52,\,6.18]$ & $+3.47\,[0.67,\,6.59]$ \\
2 & Direction-aware 96 & $+0.00\,[0.00,\,0.03]$ & $+0.00\,[0.00,\,0.00]$ & $+4.21\,[2.51,\,6.17]$ & $+3.47\,[0.67,\,6.59]$ \\
\bottomrule
\end{tabular}
\end{table*}


\clearpage
\subsection*{Absolute sparse-head results and anatomy groups}
\begin{table}[H]
\centering
\footnotesize\setlength{\tabcolsep}{4pt}
\caption{Absolute CPM and sensitivity (\%) at the seven FP/scan rates for every sparse readout head.}
\label{tab:attention_absolute}
\begin{tabular}{@{}llrrrrrrrr@{}}
\toprule
Seed & Model & CPM & 0.125 & 0.25 & 0.5 & 1 & 2 & 4 & 8 \\
\midrule
1 & Ball & 83.58 & 70.07 & 76.14 & 81.78 & 85.68 & 88.07 & 90.67 & 92.62 \\
1 & Compact (initial) & 84.44 & 70.07 & 76.57 & 82.43 & 87.20 & 89.15 & 91.76 & 93.93 \\
1 & Compact refit & 85.06 & 71.37 & 77.66 & 82.21 & 88.07 & 90.02 & 91.76 & 94.36 \\
1 & Direction-blind 48 & 84.38 & 71.15 & 75.05 & 82.43 & 86.98 & 89.59 & 91.54 & 93.93 \\
1 & Direction-aware 48 & 84.57 & 69.85 & 77.22 & 81.56 & 86.33 & 90.02 & 92.41 & 94.58 \\
1 & Direction-blind 96 & 84.54 & 70.07 & 76.57 & 82.21 & 87.42 & 89.59 & 91.54 & 94.36 \\
1 & Direction-aware 96 & 84.38 & 70.28 & 76.36 & 82.21 & 86.55 & 89.59 & 91.32 & 94.36 \\
1 & Band & 87.60 & 73.97 & 80.69 & 85.47 & 89.80 & 92.41 & 94.58 & 96.31 \\
1 & Full & 89.37 & 79.18 & 82.65 & 87.64 & 90.67 & 93.71 & 95.23 & 96.53 \\
\midrule
2 & Ball & 82.83 & 65.51 & 74.40 & 80.04 & 85.68 & 89.59 & 91.54 & 93.06 \\
2 & Compact (initial) & 82.96 & 65.08 & 74.40 & 80.48 & 85.90 & 89.37 & 91.54 & 93.93 \\
2 & Compact refit & 82.93 & 65.51 & 74.40 & 80.26 & 85.47 & 89.80 & 91.54 & 93.49 \\
2 & Direction-blind 48 & 82.83 & 65.51 & 74.40 & 80.04 & 85.68 & 89.59 & 91.54 & 93.06 \\
2 & Direction-aware 48 & 82.62 & 65.29 & 74.19 & 80.04 & 85.90 & 89.15 & 90.89 & 92.84 \\
2 & Direction-blind 96 & 82.83 & 65.51 & 74.40 & 80.04 & 85.68 & 89.59 & 91.54 & 93.06 \\
2 & Direction-aware 96 & 82.83 & 65.51 & 74.40 & 80.04 & 85.68 & 89.59 & 91.54 & 93.06 \\
2 & Band & 87.05 & 72.67 & 79.61 & 85.25 & 89.15 & 92.62 & 94.36 & 95.66 \\
2 & Full & 88.63 & 77.01 & 80.69 & 87.20 & 90.46 & 93.93 & 95.01 & 96.10 \\
\bottomrule
\end{tabular}
\end{table}

\begin{table}[H]
\centering
\small
\caption{Anatomy groups over all matching proposals, defined before scores were linked. Pairs count strict mark--lesion matches, so one lesion can contribute several pairs.}
\label{tab:anatomy_counts}
\begin{tabular}{@{}lrrrr@{}}
\toprule
Group & Lesions & With proposal & Pairs & Median diameter (mm) \\
\midrule
Does not reach band & 377 & 377 & 377 & 5.75 \\
Extends into band & 72 & 72 & 221 & 17.36 \\
No matching proposal & 6 & 0 & 0 & 4.25 \\
No contour link & 6 & 6 & 6 & 5.95 \\
\bottomrule
\end{tabular}
\end{table}

\begin{table*}[!t]
\centering
\footnotesize\setlength{\tabcolsep}{4pt}
\caption{Band-minus-ball sensitivity (points) by anatomy group at each FP/scan rate, with the 100-mark cap. The last row of each seed is the difference of the two group gains. Each cell gives the estimate above its pointwise 95\% interval.}
\label{tab:anatomy_capped}
\begin{tabular}{@{}llrrrrrrr@{}}
\toprule
Seed & Group & 0.125 & 0.25 & 0.5 & 1 & 2 & 4 & 8 \\
\midrule
1 & Does not reach band & \raisebox{-.4\height}{\shortstack{$+0.0$\\$[-3.9,\,4.7]$}} & \raisebox{-.4\height}{\shortstack{$+1.1$\\$[-3.6,\,4.5]$}} & \raisebox{-.4\height}{\shortstack{$+0.5$\\$[-3.0,\,3.8]$}} & \raisebox{-.4\height}{\shortstack{$+2.1$\\$[-0.5,\,4.8]$}} & \raisebox{-.4\height}{\shortstack{$+2.9$\\$[1.1,\,5.0]$}} & \raisebox{-.4\height}{\shortstack{$+2.1$\\$[0.5,\,4.2]$}} & \raisebox{-.4\height}{\shortstack{$+3.4$\\$[1.3,\,5.8]$}} \\
1 & Extends into band & \raisebox{-.4\height}{\shortstack{$+25.0$\\$[13.2,\,36.5]$}} & \raisebox{-.4\height}{\shortstack{$+23.6$\\$[10.1,\,33.3]$}} & \raisebox{-.4\height}{\shortstack{$+20.8$\\$[9.5,\,30.7]$}} & \raisebox{-.4\height}{\shortstack{$+15.3$\\$[5.8,\,25.0]$}} & \raisebox{-.4\height}{\shortstack{$+12.5$\\$[4.7,\,22.6]$}} & \raisebox{-.4\height}{\shortstack{$+13.9$\\$[6.2,\,22.1]$}} & \raisebox{-.4\height}{\shortstack{$+5.6$\\$[-1.4,\,12.9]$}} \\
1 & No matching proposal & \raisebox{-.4\height}{\shortstack{$+0.0$\\$[0.0,\,0.0]$}} & \raisebox{-.4\height}{\shortstack{$+0.0$\\$[0.0,\,0.0]$}} & \raisebox{-.4\height}{\shortstack{$+0.0$\\$[0.0,\,0.0]$}} & \raisebox{-.4\height}{\shortstack{$+0.0$\\$[0.0,\,0.0]$}} & \raisebox{-.4\height}{\shortstack{$+0.0$\\$[0.0,\,0.0]$}} & \raisebox{-.4\height}{\shortstack{$+0.0$\\$[0.0,\,0.0]$}} & \raisebox{-.4\height}{\shortstack{$+0.0$\\$[0.0,\,0.0]$}} \\
1 & No contour link & \raisebox{-.4\height}{\shortstack{$+0.0$\\$[0.0,\,0.0]$}} & \raisebox{-.4\height}{\shortstack{$+0.0$\\$[0.0,\,0.0]$}} & \raisebox{-.4\height}{\shortstack{$+0.0$\\$[0.0,\,0.0]$}} & \raisebox{-.4\height}{\shortstack{$+0.0$\\$[0.0,\,0.0]$}} & \raisebox{-.4\height}{\shortstack{$+0.0$\\$[0.0,\,0.0]$}} & \raisebox{-.4\height}{\shortstack{$+0.0$\\$[0.0,\,0.0]$}} & \raisebox{-.4\height}{\shortstack{$+0.0$\\$[0.0,\,0.0]$}} \\
1 & Extends $-$ does not reach & \raisebox{-.4\height}{\shortstack{$+25.0$\\$[12.5,\,37.1]$}} & \raisebox{-.4\height}{\shortstack{$+22.6$\\$[9.6,\,33.0]$}} & \raisebox{-.4\height}{\shortstack{$+20.3$\\$[8.5,\,31.0]$}} & \raisebox{-.4\height}{\shortstack{$+13.2$\\$[3.7,\,22.9]$}} & \raisebox{-.4\height}{\shortstack{$+9.6$\\$[1.4,\,20.1]$}} & \raisebox{-.4\height}{\shortstack{$+11.8$\\$[3.9,\,20.2]$}} & \raisebox{-.4\height}{\shortstack{$+2.1$\\$[-5.0,\,10.1]$}} \\
\midrule
2 & Does not reach band & \raisebox{-.4\height}{\shortstack{$+5.6$\\$[-1.0,\,11.9]$}} & \raisebox{-.4\height}{\shortstack{$+2.9$\\$[0.3,\,7.3]$}} & \raisebox{-.4\height}{\shortstack{$+3.2$\\$[0.6,\,6.6]$}} & \raisebox{-.4\height}{\shortstack{$+1.1$\\$[-1.7,\,4.3]$}} & \raisebox{-.4\height}{\shortstack{$+1.6$\\$[-0.5,\,3.6]$}} & \raisebox{-.4\height}{\shortstack{$+1.6$\\$[-0.6,\,3.8]$}} & \raisebox{-.4\height}{\shortstack{$+1.9$\\$[0.3,\,3.6]$}} \\
2 & Extends into band & \raisebox{-.4\height}{\shortstack{$+16.7$\\$[9.1,\,27.5]$}} & \raisebox{-.4\height}{\shortstack{$+18.1$\\$[6.5,\,28.9]$}} & \raisebox{-.4\height}{\shortstack{$+16.7$\\$[5.5,\,28.2]$}} & \raisebox{-.4\height}{\shortstack{$+16.7$\\$[6.8,\,26.1]$}} & \raisebox{-.4\height}{\shortstack{$+11.1$\\$[3.6,\,19.4]$}} & \raisebox{-.4\height}{\shortstack{$+9.7$\\$[2.7,\,19.0]$}} & \raisebox{-.4\height}{\shortstack{$+6.9$\\$[0.0,\,14.5]$}} \\
2 & No matching proposal & \raisebox{-.4\height}{\shortstack{$+0.0$\\$[0.0,\,0.0]$}} & \raisebox{-.4\height}{\shortstack{$+0.0$\\$[0.0,\,0.0]$}} & \raisebox{-.4\height}{\shortstack{$+0.0$\\$[0.0,\,0.0]$}} & \raisebox{-.4\height}{\shortstack{$+0.0$\\$[0.0,\,0.0]$}} & \raisebox{-.4\height}{\shortstack{$+0.0$\\$[0.0,\,0.0]$}} & \raisebox{-.4\height}{\shortstack{$+0.0$\\$[0.0,\,0.0]$}} & \raisebox{-.4\height}{\shortstack{$+0.0$\\$[0.0,\,0.0]$}} \\
2 & No contour link & \raisebox{-.4\height}{\shortstack{$+0.0$\\$[0.0,\,0.0]$}} & \raisebox{-.4\height}{\shortstack{$+0.0$\\$[0.0,\,0.0]$}} & \raisebox{-.4\height}{\shortstack{$+0.0$\\$[0.0,\,0.0]$}} & \raisebox{-.4\height}{\shortstack{$+0.0$\\$[0.0,\,0.0]$}} & \raisebox{-.4\height}{\shortstack{$+0.0$\\$[0.0,\,0.0]$}} & \raisebox{-.4\height}{\shortstack{$+0.0$\\$[0.0,\,0.0]$}} & \raisebox{-.4\height}{\shortstack{$+0.0$\\$[0.0,\,0.0]$}} \\
2 & Extends $-$ does not reach & \raisebox{-.4\height}{\shortstack{$+11.1$\\$[2.7,\,23.2]$}} & \raisebox{-.4\height}{\shortstack{$+15.1$\\$[2.2,\,26.1]$}} & \raisebox{-.4\height}{\shortstack{$+13.5$\\$[1.9,\,25.3]$}} & \raisebox{-.4\height}{\shortstack{$+15.6$\\$[4.7,\,24.9]$}} & \raisebox{-.4\height}{\shortstack{$+9.5$\\$[1.7,\,18.0]$}} & \raisebox{-.4\height}{\shortstack{$+8.1$\\$[0.8,\,17.8]$}} & \raisebox{-.4\height}{\shortstack{$+5.1$\\$[-2.9,\,12.8]$}} \\
\bottomrule
\end{tabular}
\end{table*}


\clearpage
\subsection*{Anatomy without the mark cap, per-lesion decisions and earlier comparisons}
\begin{table*}[!t]
\centering
\footnotesize\setlength{\tabcolsep}{4pt}
\caption{As Tab.~\ref{tab:anatomy_capped}, without the 100-mark cap.}
\label{tab:anatomy_uncapped}
\begin{tabular}{@{}llrrrrrrr@{}}
\toprule
Seed & Group & 0.125 & 0.25 & 0.5 & 1 & 2 & 4 & 8 \\
\midrule
1 & Does not reach band & \raisebox{-.4\height}{\shortstack{$+0.0$\\$[-3.9,\,4.7]$}} & \raisebox{-.4\height}{\shortstack{$+1.1$\\$[-3.6,\,4.5]$}} & \raisebox{-.4\height}{\shortstack{$+0.5$\\$[-3.0,\,3.8]$}} & \raisebox{-.4\height}{\shortstack{$+2.1$\\$[-0.5,\,4.8]$}} & \raisebox{-.4\height}{\shortstack{$+2.9$\\$[1.1,\,5.0]$}} & \raisebox{-.4\height}{\shortstack{$+2.1$\\$[0.5,\,4.2]$}} & \raisebox{-.4\height}{\shortstack{$+3.4$\\$[1.3,\,5.8]$}} \\
1 & Extends into band & \raisebox{-.4\height}{\shortstack{$+25.0$\\$[13.2,\,36.5]$}} & \raisebox{-.4\height}{\shortstack{$+23.6$\\$[10.1,\,33.3]$}} & \raisebox{-.4\height}{\shortstack{$+20.8$\\$[9.5,\,30.7]$}} & \raisebox{-.4\height}{\shortstack{$+15.3$\\$[5.8,\,25.0]$}} & \raisebox{-.4\height}{\shortstack{$+12.5$\\$[4.7,\,22.6]$}} & \raisebox{-.4\height}{\shortstack{$+13.9$\\$[6.2,\,22.1]$}} & \raisebox{-.4\height}{\shortstack{$+5.6$\\$[-1.4,\,12.9]$}} \\
1 & No matching proposal & \raisebox{-.4\height}{\shortstack{$+0.0$\\$[0.0,\,0.0]$}} & \raisebox{-.4\height}{\shortstack{$+0.0$\\$[0.0,\,0.0]$}} & \raisebox{-.4\height}{\shortstack{$+0.0$\\$[0.0,\,0.0]$}} & \raisebox{-.4\height}{\shortstack{$+0.0$\\$[0.0,\,0.0]$}} & \raisebox{-.4\height}{\shortstack{$+0.0$\\$[0.0,\,0.0]$}} & \raisebox{-.4\height}{\shortstack{$+0.0$\\$[0.0,\,0.0]$}} & \raisebox{-.4\height}{\shortstack{$+0.0$\\$[0.0,\,0.0]$}} \\
1 & No contour link & \raisebox{-.4\height}{\shortstack{$+0.0$\\$[0.0,\,0.0]$}} & \raisebox{-.4\height}{\shortstack{$+0.0$\\$[0.0,\,0.0]$}} & \raisebox{-.4\height}{\shortstack{$+0.0$\\$[0.0,\,0.0]$}} & \raisebox{-.4\height}{\shortstack{$+0.0$\\$[0.0,\,0.0]$}} & \raisebox{-.4\height}{\shortstack{$+0.0$\\$[0.0,\,0.0]$}} & \raisebox{-.4\height}{\shortstack{$+0.0$\\$[0.0,\,0.0]$}} & \raisebox{-.4\height}{\shortstack{$+0.0$\\$[0.0,\,0.0]$}} \\
1 & Extends $-$ does not reach & \raisebox{-.4\height}{\shortstack{$+25.0$\\$[12.5,\,37.1]$}} & \raisebox{-.4\height}{\shortstack{$+22.6$\\$[9.6,\,33.0]$}} & \raisebox{-.4\height}{\shortstack{$+20.3$\\$[8.5,\,31.0]$}} & \raisebox{-.4\height}{\shortstack{$+13.2$\\$[3.7,\,22.9]$}} & \raisebox{-.4\height}{\shortstack{$+9.6$\\$[1.4,\,20.1]$}} & \raisebox{-.4\height}{\shortstack{$+11.8$\\$[3.9,\,20.2]$}} & \raisebox{-.4\height}{\shortstack{$+2.1$\\$[-5.0,\,10.1]$}} \\
\midrule
2 & Does not reach band & \raisebox{-.4\height}{\shortstack{$+5.6$\\$[-1.0,\,11.9]$}} & \raisebox{-.4\height}{\shortstack{$+2.9$\\$[0.3,\,7.3]$}} & \raisebox{-.4\height}{\shortstack{$+3.2$\\$[0.6,\,6.6]$}} & \raisebox{-.4\height}{\shortstack{$+1.1$\\$[-1.7,\,4.3]$}} & \raisebox{-.4\height}{\shortstack{$+1.6$\\$[-0.5,\,3.6]$}} & \raisebox{-.4\height}{\shortstack{$+1.6$\\$[-0.6,\,3.8]$}} & \raisebox{-.4\height}{\shortstack{$+1.9$\\$[0.3,\,3.6]$}} \\
2 & Extends into band & \raisebox{-.4\height}{\shortstack{$+16.7$\\$[9.1,\,27.5]$}} & \raisebox{-.4\height}{\shortstack{$+18.1$\\$[6.5,\,28.9]$}} & \raisebox{-.4\height}{\shortstack{$+16.7$\\$[5.5,\,28.2]$}} & \raisebox{-.4\height}{\shortstack{$+16.7$\\$[6.8,\,26.1]$}} & \raisebox{-.4\height}{\shortstack{$+11.1$\\$[3.6,\,19.4]$}} & \raisebox{-.4\height}{\shortstack{$+9.7$\\$[2.7,\,19.0]$}} & \raisebox{-.4\height}{\shortstack{$+6.9$\\$[0.0,\,14.8]$}} \\
2 & No matching proposal & \raisebox{-.4\height}{\shortstack{$+0.0$\\$[0.0,\,0.0]$}} & \raisebox{-.4\height}{\shortstack{$+0.0$\\$[0.0,\,0.0]$}} & \raisebox{-.4\height}{\shortstack{$+0.0$\\$[0.0,\,0.0]$}} & \raisebox{-.4\height}{\shortstack{$+0.0$\\$[0.0,\,0.0]$}} & \raisebox{-.4\height}{\shortstack{$+0.0$\\$[0.0,\,0.0]$}} & \raisebox{-.4\height}{\shortstack{$+0.0$\\$[0.0,\,0.0]$}} & \raisebox{-.4\height}{\shortstack{$+0.0$\\$[0.0,\,0.0]$}} \\
2 & No contour link & \raisebox{-.4\height}{\shortstack{$+0.0$\\$[0.0,\,0.0]$}} & \raisebox{-.4\height}{\shortstack{$+0.0$\\$[0.0,\,0.0]$}} & \raisebox{-.4\height}{\shortstack{$+0.0$\\$[0.0,\,0.0]$}} & \raisebox{-.4\height}{\shortstack{$+0.0$\\$[0.0,\,0.0]$}} & \raisebox{-.4\height}{\shortstack{$+0.0$\\$[0.0,\,0.0]$}} & \raisebox{-.4\height}{\shortstack{$+0.0$\\$[0.0,\,0.0]$}} & \raisebox{-.4\height}{\shortstack{$+0.0$\\$[0.0,\,0.0]$}} \\
2 & Extends $-$ does not reach & \raisebox{-.4\height}{\shortstack{$+11.1$\\$[2.7,\,23.2]$}} & \raisebox{-.4\height}{\shortstack{$+15.1$\\$[2.2,\,26.1]$}} & \raisebox{-.4\height}{\shortstack{$+13.5$\\$[1.9,\,25.3]$}} & \raisebox{-.4\height}{\shortstack{$+15.6$\\$[4.7,\,24.9]$}} & \raisebox{-.4\height}{\shortstack{$+9.5$\\$[1.7,\,18.0]$}} & \raisebox{-.4\height}{\shortstack{$+8.1$\\$[0.8,\,17.8]$}} & \raisebox{-.4\height}{\shortstack{$+5.1$\\$[-2.4,\,13.1]$}} \\
\bottomrule
\end{tabular}
\end{table*}

\begin{table}[H]
\centering
\footnotesize\setlength{\tabcolsep}{4pt}
\caption{Per-lesion decisions of seed 1 and seed 2 at the empirical and transferred cutoffs (descriptive). The last two columns count lesions whose decision mark before or after restoring the band has a band that reaches the majority-reader contour.}
\label{tab:decision}
\begin{tabular}{@{}lllrrr@{}}
\toprule
Seed & Cutoff & Status & Lesions & Before overlap & After overlap \\
\midrule
1 & Empirical & Retained & 386 & 44 & 45 \\
1 & Empirical & Gained & 28 & 0 & 12 \\
1 & Empirical & Lost & 9 & 2 & 0 \\
1 & Empirical & Missed & 38 & 0 & 0 \\
\midrule
1 & Transfer & Retained & 386 & 44 & 45 \\
1 & Transfer & Gained & 28 & 0 & 12 \\
1 & Transfer & Lost & 11 & 3 & 0 \\
1 & Transfer & Missed & 36 & 0 & 0 \\
\midrule
2 & Empirical & Retained & 383 & 45 & 46 \\
2 & Empirical & Gained & 28 & 0 & 12 \\
2 & Empirical & Lost & 12 & 1 & 0 \\
2 & Empirical & Missed & 38 & 0 & 0 \\
\midrule
2 & Transfer & Retained & 384 & 45 & 46 \\
2 & Transfer & Gained & 25 & 0 & 11 \\
2 & Transfer & Lost & 14 & 2 & 0 \\
2 & Transfer & Missed & 38 & 0 & 0 \\
\bottomrule
\end{tabular}
\end{table}

\begin{table}[H]
\centering
\small
\caption{Earlier comparisons with classical half-axis readouts and visibility on the same benchmark. Pointwise 95\% patient-level intervals.}
\label{tab:historical}
\begin{tabular}{@{}lrr@{}}
\toprule
Contrast & $\Delta$ CPM & $\Delta S_1$ \\
\midrule
Fixed half-axes $-$ raw Hessian & $+3.32\,[2.52,\,4.33]$ & $+4.99\,[2.52,\,7.28]$ \\
Fixed half-axes $-$ shuffled & $+1.60\,[0.48,\,3.15]$ & $+2.79\,[0.09,\,5.23]$ \\
Tuned half-axes $-$ shuffled & $+2.08\,[0.98,\,3.19]$ & $+4.04\,[0.91,\,5.77]$ \\
Tuned half-axes $-$ grey max. & $-0.71\,[-2.98,\,1.46]$ & $+0.00\,[-4.05,\,3.02]$ \\
Visibility $\times$ half-axis interaction & $-0.03\,[-0.76,\,0.46]$ & $-0.22\,[-1.51,\,1.04]$ \\
\bottomrule
\end{tabular}
\end{table}

% Transcribed from the authors' BLUE_REVIEW_33391_v2 supplement (exploratory size/displacement audit).
\begin{table}[h]
\centering
\small
\caption{Exploratory size/displacement audit at calibration-fixed thresholds: band-minus-ball sensitivity gain (points) per cell of reference diameter and minimum matching-proposal displacement. O: the lesion extends into the band. Cells with fewer than 10 lesions are descriptive only.}
\label{tab:size_displacement}
\begin{tabular}{@{}llrrrr@{}}
\toprule
Size (mm) & Offset (mm) & O & $n$ & Seed 1 & Seed 2 \\
\midrule
$<6$ & $<1.5$ & 0 & 198 & +2.02 & +0.00 \\
$<6$ & $<1.5$ & 1 & 0 & -- & -- \\
$<6$ & 1.5--$<3$ & 0 & 9 & +0.00 & $-$11.11 \\
$<6$ & 1.5--$<3$ & 1 & 0 & -- & -- \\
$<6$ & $\ge3$ & 0 & 0 & -- & -- \\
$<6$ & $\ge3$ & 1 & 0 & -- & -- \\
6--$<10$ & $<1.5$ & 0 & 117 & +0.85 & +1.71 \\
6--$<10$ & $<1.5$ & 1 & 0 & -- & -- \\
6--$<10$ & 1.5--$<3$ & 0 & 17 & +17.65 & +0.00 \\
6--$<10$ & 1.5--$<3$ & 1 & 1 & +0.00 & +0.00 \\
6--$<10$ & $\ge3$ & 0 & 1 & $-$100.00 & $-$100.00 \\
6--$<10$ & $\ge3$ & 1 & 1 & +0.00 & +0.00 \\
$\ge10$ & $<1.5$ & 0 & 24 & +0.00 & +0.00 \\
$\ge10$ & $<1.5$ & 1 & 36 & +13.89 & +5.56 \\
$\ge10$ & 1.5--$<3$ & 0 & 11 & +0.00 & +9.09 \\
$\ge10$ & 1.5--$<3$ & 1 & 30 & +13.33 & +20.00 \\
$\ge10$ & $\ge3$ & 0 & 0 & -- & -- \\
$\ge10$ & $\ge3$ & 1 & 4 & +25.00 & +50.00 \\
\bottomrule
\end{tabular}
\end{table}

